\documentclass[12pt]{article}

\usepackage[a4paper, margin=2.5cm]{geometry}
\usepackage{amsmath}
\usepackage{amssymb}
\usepackage{amsthm}
\usepackage[colorlinks=true, linkcolor=blue, citecolor=blue, urlcolor=blue]{hyperref}
\usepackage{booktabs}
\usepackage{natbib}

\newtheorem{theorem}{Theorem}
\newtheorem{proposition}[theorem]{Proposition}
\newtheorem{corollary}[theorem]{Corollary}
\newtheorem{definition}[theorem]{Definition}
\newtheorem{remark}[theorem]{Remark}
\newtheorem{example}[theorem]{Example}

\newcommand{\Fcal}{\mathcal{F}}
\newcommand{\Ncl}{\mathcal{N}_{\mathrm{cl}}}
\newcommand{\VW}{V_{\mathrm{W}}}
\newcommand{\HW}{H_{\mathrm{W}}}
\newcommand{\Sadm}{\mathcal{S}_{\mathrm{adm}}}
\newcommand{\Rst}{R_{\mathrm{st}}}
\newcommand{\dUrel}{\delta U_{\mathrm{rel}}}

\title{The Closure Residual and Nonlinear Quantum Structure\\[6pt]
\large Equilibrium Quantum Limit and Modulus-Curvature Nonlinearity}
\author{%
  Lee Smart\\[4pt]
  \textit{Institute of Vibrational Field Dynamics}\\[2pt]
  \texttt{contact@vibrationalfielddynamics.org}\\[2pt]
  \texttt{@vfd\_org}%
}
\date{April 2026}

\begin{document}

\maketitle

\noindent\textbf{Status:} Preprint (not peer-reviewed)

\begin{abstract}
Paper~XVIII derived an exact nonlinear closure-Nelson wave equation and showed that, at the Paper~XIV equilibrium and to leading order near equilibrium, the dynamics reduces to the Schr\"odinger equation with Witten Hamiltonian $\HW$. Paper~XIX isolated the remaining nonlinear term as the closure residual. The present paper places that residual into nonlinear quantum structure.

Writing the dynamics in the equilibrium-centred form
\[
i\partial_t\Psi = \HW\Psi + \Ncl[\Psi],
\]
the nonlinear operator $\Ncl$ is shown to be a local modulus-curvature term with equilibrium subtraction, fixed entirely by the closure geometry. Its structural properties are derived: phase covariance, amplitude homogeneity, locality away from nodal sets, exact vanishing on the equilibrium amplitude, and perturbative smallness near equilibrium. These results identify the linear Schr\"odinger equation not as the full closure dynamics, but as the distinguished equilibrium tangent limit of an exact nonlinear wave equation.

The closure nonlinearity is then compared structurally with known classes of nonlinear Schr\"odinger equations, including modulus-dependent curvature corrections. The claim is not exact identification with any existing family, but a precise mathematical placement: the closure residual defines a derived modulus-curvature nonlinear Schr\"odinger structure whose coefficient and equilibrium subtraction are fixed by the closure geometry.

This paper does not assign physical interpretation to the nonlinear residual. It establishes that the closure framework yields a mathematically well-defined nonlinear quantum extension with an exact equilibrium quantum limit.
\end{abstract}

%==================================================================
\section{Introduction}\label{sec:intro}
%==================================================================

The closure programme has now constructed a layered route from classical gradient-flow dynamics to quantum-type evolution:
\begin{itemize}
    \item Paper~XVI~\citep{SmartXVI}: closure evolution equation and structural gap to Schr\"odinger,
    \item Paper~XVII~\citep{SmartXVII}: dissipative obstruction for single-generator approaches,
    \item Paper~XVIII~\citep{SmartXVIII}: Nelson pairing and near-equilibrium Schr\"odinger recovery,
    \item Paper~XIX~\citep{SmartXIX}: isolation and classification of the nonlinear residual.
\end{itemize}

Paper~XIX established that the residual is a robust, non-trivial, probability-conserving nonlinear remainder that cannot be removed by standard gauge transformations. The purpose of the present paper is not to re-derive that remainder, but to place it within the broader structure of nonlinear Schr\"odinger dynamics and identify the linear quantum regime as its equilibrium tangent limit.

\subsection*{Purpose}

This paper:
\begin{enumerate}
    \item defines the closure nonlinear operator $\Ncl$ and establishes its structural properties,
    \item proves that the linear Schr\"odinger equation is the equilibrium tangent limit of the full nonlinear dynamics,
    \item decomposes the nonlinearity into a universal modulus-curvature term and an equilibrium-fixed subtraction,
    \item compares the closure nonlinearity with known nonlinear Schr\"odinger classes.
\end{enumerate}

It does not claim experimental predictions or exact identification with existing nonlinear QM models. It is a structural classification paper.

%==================================================================
\section{The Closure Nonlinear Operator and Equilibrium-Centred Form}\label{sec:operator}
%==================================================================

\subsection{Definition}

Starting from the Paper~XIX normal form, define the \textit{closure nonlinear operator}:

\begin{definition}[Closure nonlinear operator]\label{def:Ncl}
\begin{equation}\label{eq:Ncl}
    \Ncl[\Psi] = \dUrel[\Psi]\,\Psi = \sigma^2\left(\frac{\Delta|\Psi|}{|\Psi|} - \frac{\Delta\Rst}{\Rst}\right)\Psi
\end{equation}
where $\Rst = \sqrt{\rho_{\mathrm{st}}} \propto e^{-\Fcal/\sigma^2}$ and $\Delta\Rst/\Rst = 2\VW/\sigma^2$.
\end{definition}

The full closure-Nelson dynamics is then:
\begin{equation}\label{eq:full}
    i\,\partial_t\Psi = \HW\Psi + \Ncl[\Psi]
\end{equation}

\subsection{Decomposition}

The operator splits into two pieces:
\begin{equation}\label{eq:decomposition}
    \Ncl[\Psi] = \underbrace{\sigma^2\frac{\Delta|\Psi|}{|\Psi|}\,\Psi}_{\text{universal modulus-curvature term}} - \underbrace{2\VW\,\Psi}_{\text{equilibrium subtraction}}
\end{equation}

The first term depends only on the modulus profile of $\Psi$ and is the same functional form that appears in the quantum potential of the Madelung formulation. The second term is entirely fixed by the closure equilibrium geometry: it subtracts the equilibrium curvature so that $\Ncl$ vanishes on $\Psi_{\mathrm{st}}$.

\begin{remark}
The equilibrium subtraction is not a free parameter. It is determined by $\Fcal$ and $\sigma$ through $\Delta\Rst/\Rst = 2\VW/\sigma^2$. This is the key structural difference between the closure nonlinearity and a generic nonlinear Schr\"odinger term: the subtraction is \textit{derived from the closure landscape}, not introduced phenomenologically.
\end{remark}

%==================================================================
\section{Structural Properties}\label{sec:structure}
%==================================================================

\begin{proposition}[Phase covariance]\label{prop:covariance}
For any constant $\theta \in \mathbb{R}$:
\begin{equation}
    \Ncl[e^{i\theta}\Psi] = e^{i\theta}\Ncl[\Psi]
\end{equation}
That is, $\Ncl$ is covariant under global phase rotations.
\end{proposition}

\begin{proof}
$|e^{i\theta}\Psi| = |\Psi|$, so $\dUrel[e^{i\theta}\Psi] = \dUrel[\Psi]$, and $\Ncl[e^{i\theta}\Psi] = \dUrel[\Psi]\cdot e^{i\theta}\Psi = e^{i\theta}\Ncl[\Psi]$.
\end{proof}

\begin{proposition}[Amplitude homogeneity]\label{prop:homogeneity}
For any constant $\lambda > 0$:
\begin{equation}
    \Ncl[\lambda\Psi] = \lambda\,\Ncl[\Psi]
\end{equation}
\end{proposition}

\begin{proof}
$|\lambda\Psi| = \lambda|\Psi|$, so $\Delta|\lambda\Psi|/|\lambda\Psi| = \Delta|\Psi|/|\Psi|$. Then $\dUrel[\lambda\Psi] = \dUrel[\Psi]$, and $\Ncl[\lambda\Psi] = \dUrel[\Psi]\cdot\lambda\Psi = \lambda\Ncl[\Psi]$.
\end{proof}

\begin{corollary}[Combined scaling]
$\Ncl[\lambda e^{i\theta}\Psi] = \lambda e^{i\theta}\Ncl[\Psi]$ for $\lambda > 0$, $\theta \in \mathbb{R}$. The operator is homogeneous of degree one and phase-covariant.
\end{corollary}

\noindent\textit{Proof.} Immediate from Propositions~\ref{prop:covariance} and~\ref{prop:homogeneity}. \hfill$\square$

\begin{proposition}[Locality]\label{prop:locality}
On any open set where $|\Psi| > 0$, $\Ncl[\Psi](\psi)$ depends only on $|\Psi|$, $\nabla|\Psi|$, $\Delta|\Psi|$, $\Psi$, and the equilibrium data $\Rst$ evaluated at $\psi$. It is a local nonlinear operator.
\end{proposition}

\begin{proof}
Immediate from the definition~\eqref{eq:Ncl}: $\Delta|\Psi|/|\Psi|$ involves only the values and first two spatial derivatives of $|\Psi|$ at the point $\psi$, and $\Delta\Rst/\Rst$ is a fixed function of $\psi$.
\end{proof}

\begin{proposition}[Equilibrium vanishing]\label{prop:vanishing}
\begin{equation}
    \Ncl[\Psi_{\mathrm{st}}] = 0
\end{equation}
The equilibrium state lies in the kernel of the nonlinear operator. On the equilibrium amplitude sector, the dynamics is exactly linear Schr\"odinger.
\end{proposition}

\begin{proof}
Direct: $|\Psi_{\mathrm{st}}| = \Rst$, so $\Delta|\Psi_{\mathrm{st}}|/|\Psi_{\mathrm{st}}| = \Delta\Rst/\Rst$ and the difference in~\eqref{eq:Ncl} vanishes.
\end{proof}

\begin{proposition}[Modulus-only dependence]\label{prop:modulus}
$\Ncl[\Psi]$ depends on the phase of $\Psi$ only through the overall factor $\Psi$ in~\eqref{eq:Ncl}. The residual potential $\dUrel[\Psi]$ is a functional of $|\Psi|$ alone.
\end{proposition}

\begin{proof}
The quantity $\Delta|\Psi|/|\Psi|$ depends only on the modulus and its derivatives, not on the phase of $\Psi$. The equilibrium term $\Delta\Rst/\Rst$ is phase-independent. Thus $\dUrel[\Psi]$ is modulus-only; the factor $\Psi$ in $\Ncl[\Psi] = \dUrel[\Psi]\Psi$ carries the phase.
\end{proof}

%==================================================================
\section{The Equilibrium Quantum Limit}\label{sec:limit}
%==================================================================

The linear Schr\"odinger equation is not imposed on the closure framework; it emerges as the distinguished tangent dynamics at equilibrium.

\begin{proposition}[Equilibrium tangent limit]\label{prop:tangent}
Let $\Psi(t) = \Psi_{\mathrm{st}}(1 + \varepsilon\,\eta(\psi,t))$ with $\varepsilon \ll 1$. Then the Fr\'echet linearisation of the full nonlinear flow~\eqref{eq:full} at the equilibrium state $\Psi_{\mathrm{st}}$ is the Witten-Hamiltonian Schr\"odinger generator: to leading order,
\begin{equation}
    i\,\partial_t(\Psi_{\mathrm{st}}\,\eta) = \HW(\Psi_{\mathrm{st}}\,\eta) + O(\varepsilon)
\end{equation}
That is, infinitesimal perturbations around equilibrium evolve under the linear Schr\"odinger equation $i\partial_t\delta\Psi = \HW\,\delta\Psi$.
\end{proposition}

\begin{proof}
Since $\Ncl[\Psi_{\mathrm{st}}] = 0$ (Proposition~\ref{prop:vanishing}) and $\dUrel$ is continuous in $|\Psi|$, the residual potential is Fr\'echet-small under $\Psi = \Psi_{\mathrm{st}}(1 + \varepsilon\eta)$: $\dUrel[\Psi] = O(\varepsilon)$ (Paper~XIX, Proposition~4). The nonlinear contribution is of higher perturbative order than the linear Witten term, so the Fr\'echet linearisation of the full flow is generated by $\HW$.
\end{proof}

\begin{remark}[Interpretation]
The Witten-Hamiltonian Schr\"odinger equation $i\partial_t\Psi = \HW\Psi$ is not the full closure dynamics. It is the \textit{tangent linearisation} of the full nonlinear dynamics at the equilibrium fixed point $\Psi_{\mathrm{st}}$. Linear quantum mechanics is the equilibrium sector of the closure theory, in the same sense that linear stability analysis is the equilibrium sector of a nonlinear dynamical system.

This is a structural result: the linear quantum regime is selected by the equilibrium geometry, not by postulate.
\end{remark}

%==================================================================
\section{Relation to the Madelung Quantum Potential}\label{sec:madelung}
%==================================================================

In the Madelung formulation of standard quantum mechanics, the Hamilton--Jacobi equation contains the \textit{quantum potential}:
\begin{equation}
    Q_{\mathrm{std}} = -\frac{\hbar^2}{2m}\frac{\Delta\sqrt{\rho}}{\sqrt{\rho}}
\end{equation}

The closure nonlinearity involves the same functional form:
\begin{equation}
    \Ncl[\Psi] = \sigma^2\frac{\Delta|\Psi|}{|\Psi|}\,\Psi - 2\VW\,\Psi
\end{equation}

The first term is proportional to $(\Delta\sqrt\rho/\sqrt\rho)\Psi$ with coefficient $\sigma^2$ (in Schr\"odinger-normalised form; $\sigma^4$ in the undivided Paper~XVIII form). The structural connection is:

\begin{table}[ht]
\centering
\caption{Quantum potential in standard QM vs.\ the closure nonlinearity.}
\label{tab:qp}
\begin{tabular}{@{}lll@{}}
\toprule
Feature & Standard QM & Closure framework \\
\midrule
Functional form & $\Delta\sqrt\rho/\sqrt\rho$ & $\Delta|\Psi|/|\Psi|$ (same) \\
Coefficient & $-\hbar^2/(2m)$ & $+\sigma^2$ (normalised) \\
Sign in HJ equation & Negative & Positive (reversed) \\
Equilibrium subtraction & None & $-2\sigma^2\VW$ (closure-derived) \\
Role in wave equation & Absorbed into linear kinetic term & Residual nonlinear correction \\
\bottomrule
\end{tabular}
\end{table}

\begin{remark}
In standard QM, the quantum potential is entirely absorbed into the linear Schr\"odinger kinetic term --- it is not a separate nonlinear correction. In the closure framework, the sign reversal (Paper~XVIII) prevents full absorption, leaving a residual nonlinear term. The equilibrium subtraction then ensures this residual vanishes at the closure ground state. The closure residual is therefore not merely ``another nonlinear term''; it is the non-cancelled remainder of the same modulus-curvature object that, in standard linear quantum mechanics, is hidden inside the Schr\"odinger kinetic structure.
\end{remark}

%==================================================================
\section{Relation to Nonlinear Schr\"odinger Classes}\label{sec:NLS}
%==================================================================

The closure equation~\eqref{eq:full} belongs to the broad class of nonlinear Schr\"odinger equations of the form:
\begin{equation}\label{eq:general_NLS}
    i\,\partial_t\Psi = \hat{H}_0\Psi + f(|\Psi|, \nabla|\Psi|, \Delta|\Psi|)\,\Psi
\end{equation}
where $\hat{H}_0$ is a linear Hamiltonian and $f$ is a real-valued nonlinear functional of the modulus.

\subsection{Modulus-curvature nonlinearities}

The closure nonlinearity has the specific form:
\begin{equation}
    f_{\mathrm{cl}} = \sigma^2\frac{\Delta|\Psi|}{|\Psi|} - 2\VW
\end{equation}

This is a \textit{modulus-curvature} nonlinearity: it depends on the Laplacian of the modulus (second-order spatial derivatives), not on powers of the modulus. This distinguishes it from:
\begin{itemize}
    \item \textbf{Gross--Pitaevskii} ($f = g|\Psi|^2$): power-law, no derivatives.
    \item \textbf{Logarithmic} ($f = \lambda\log|\Psi|^2$): no derivatives, but singular.
\end{itemize}

It is structurally closest to:
\begin{itemize}
    \item \textbf{Doebner--Goldin-type} nonlinearities, which include terms proportional to $\Delta|\Psi|/|\Psi|$ and arise from representations of the diffeomorphism group on quantum states.
\end{itemize}

\subsection{Comparison with generic modulus-curvature NLS}

A generic modulus-curvature nonlinear Schr\"odinger equation has the form:
\begin{equation}
    i\partial_t\Psi = \hat{H}_0\Psi + D\frac{\Delta|\Psi|}{|\Psi|}\Psi
\end{equation}
with $D$ a real coupling constant. The closure equation differs in two ways:

\begin{enumerate}
    \item \textbf{The coefficient is fixed.} In the closure framework, $D = \sigma^2$ (in Schr\"odinger-normalised form) is determined by the noise scale, not a free phenomenological parameter.

    \item \textbf{An equilibrium subtraction is present.} The closure equation subtracts $2\sigma^2\VW$, which ensures the nonlinearity vanishes at the closure equilibrium. A generic modulus-curvature NLS has no such subtraction.
\end{enumerate}

\begin{remark}[Structural placement]
The closure equation is not presented as identical to the Doebner--Goldin family or any specific nonlinear QM model. It may be placed within the same broad class of modulus-curvature nonlinear Schr\"odinger equations, but with two closure-specific features: a derived coefficient ($\sigma^2$ in normalised form) and a derived equilibrium subtraction ($2\VW$). These features are consequences of the closure geometry, not free parameters.
\end{remark}

\begin{remark}[On relation to nonlinear quantum theories]
The closure equation is not proposed as an alternative nonlinear quantum theory. Its role here is structural: to show that a nonlinear Schr\"odinger-type equation with a distinguished linear limit arises naturally from the closure construction. Whether this structure corresponds to a physical nonlinear extension of quantum mechanics is a separate question not addressed in this paper.
\end{remark}

%==================================================================
\section{What the Closure Framework Fixes}\label{sec:fixed}
%==================================================================

In a generic nonlinear Schr\"odinger theory of the form~\eqref{eq:general_NLS}, the nonlinear functional $f$ is typically a phenomenological input. The closure framework fixes several features that would otherwise be free:

\begin{table}[ht]
\centering
\caption{What the closure framework determines vs.\ what is free in generic NLS.}
\label{tab:fixed}
\begin{tabular}{@{}lll@{}}
\toprule
Feature & Generic NLS & Closure framework \\
\midrule
Linear Hamiltonian & Free choice & $\HW$ (Witten, from $\Fcal$, $\sigma$) \\
Nonlinear functional form & Free choice & $\Delta|\Psi|/|\Psi|$ (modulus curvature) \\
Nonlinear coefficient & Free ($D$) & Fixed ($\sigma^2$, normalised) \\
Equilibrium subtraction & Absent & Present ($2\VW$, from $\Fcal$) \\
Distinguished linear limit & Not automatic & Equilibrium tangent dynamics \\
Norm conservation & Must be checked & Guaranteed (real $\dUrel$) \\
\bottomrule
\end{tabular}
\end{table}

Within the forward/backward closure-pairing construction, the framework therefore produces a \textit{more constrained} nonlinear Schr\"odinger structure than a generic phenomenological model: fewer free parameters, a distinguished linear limit, and a derived equilibrium subtraction.

%==================================================================
\section{Worked Example: Quadratic Closure}\label{sec:example}
%==================================================================

\begin{example}[Quadratic closure]\label{ex:quadratic}
Let $\Fcal(\psi) = \frac{1}{2}k\psi^2$ in one dimension.

\subsection*{Closure data}

$\Rst \propto e^{-k\psi^2/(2\sigma^2)}$, $\VW = k^2\psi^2/(2\sigma^2) - k/2$, $\HW$ has spectrum $E_n = nk$.

\subsection*{Nonlinear operator on specific states}

\textbf{Ground state} $\Psi = \phi_0 \propto \Rst$:
\begin{equation}
    \Ncl[\phi_0] = 0 \qquad \text{(exact Schr\"odinger sector)}
\end{equation}

\textbf{Shifted Gaussian} $|\Psi| \propto e^{-k(\psi-a)^2/(2\sigma^2)}$ with displacement $a$:
\begin{equation}
    \dUrel = k^2(-2a\psi + a^2)
\end{equation}
The residual potential is linear in $\psi$ (plus constant), growing with displacement $a$. For small $a$, $|\dUrel| \ll |\VW|$ and the dynamics is in the Schr\"odinger regime.

\textbf{First excited state} $|\Psi| = |\phi_1| \propto |\psi|\,e^{-k\psi^2/(2\sigma^2)}$:
\begin{equation}
    \dUrel = 2\sigma^2/\psi^2 - 2k \qquad (\psi \neq 0)
\end{equation}
The residual diverges at the node $\psi = 0$ (where $|\Psi| = 0$) and is large near the origin. The excited eigenstate dynamics is in the nonlinear closure regime near nodes.

\subsection*{Regime classification}

\begin{table}[ht]
\centering
\caption{Residual strength for different states in the quadratic model.}
\label{tab:states}
\begin{tabular}{@{}llll@{}}
\toprule
State & $\dUrel$ & $\Xi$ & Regime \\
\midrule
Ground $\phi_0$ & $0$ & $0$ & Exact Schr\"odinger \\
Shifted Gaussian ($a \ll \sigma/k$) & $O(a)$ & $\ll 1$ & Linear Schr\"odinger \\
Shifted Gaussian ($a \gg \sigma/k$) & $O(k^2 a^2)$ & $\sim 1$ or $\gg 1$ & Mixed / nonlinear \\
Excited $\phi_1$ (away from node) & $O(\sigma^2 k)$ & $O(1)$ & Mixed \\
Excited $\phi_1$ (near node) & Divergent & $\gg 1$ & Nonlinear closure \\
\bottomrule
\end{tabular}
\end{table}

The quadratic example demonstrates the full regime structure: exact Schr\"odinger at the ground state, perturbative Schr\"odinger for small displacements, and genuinely nonlinear dynamics for large excitations or near nodal sets. Accordingly, statements involving excited states should be understood piecewise on nodal domains unless a regularised treatment is supplied.
\end{example}

%==================================================================
\section{Interpretation and Limits}\label{sec:interpretation}
%==================================================================

\subsection*{What is established}

The closure residual defines a specific nonlinear Schr\"odinger structure:
\begin{itemize}
    \item a modulus-curvature nonlinearity with derived coefficient $\sigma^2$ (in Schr\"odinger-normalised form),
    \item an equilibrium subtraction fixed by the closure geometry,
    \item a distinguished linear quantum limit at equilibrium,
    \item norm conservation in all regimes.
\end{itemize}

\subsection*{What is not claimed}

\begin{itemize}
    \item \textbf{No exact identification with Doebner--Goldin or any specific NLS family.} The closure equation belongs to the same broad class but has closure-specific features (equilibrium subtraction, derived coefficient) not present in generic models.
    \item \textbf{No experimental prediction.} Determining whether the nonlinear residual produces observable effects requires estimating $\sigma^2\Delta|\Psi|/|\Psi|$ in physically realistic regimes, which is not attempted here.
    \item \textbf{No claim about the physical status of the residual.} Whether it represents genuine beyond-QM physics, an effective-theory artefact, or a gauge-removable structure remains open.
\end{itemize}

\subsection*{The structural message}

The closure framework does not terminate at the linear quantum regime. It produces a nonlinear Schr\"odinger structure whose exact linear limit is the equilibrium-centred Witten-Hamiltonian regime. The nonlinearity is not a defect; it is a structural prediction of the paired-process construction, with a coefficient and subtraction term determined by the closure geometry.

Whether this prediction is physical is a question for subsequent work.

%==================================================================
\section{Discussion}\label{sec:discussion}
%==================================================================

\subsection{Exact results}

\begin{itemize}
    \item The closure nonlinear operator $\Ncl$ is defined and its decomposition into modulus-curvature and equilibrium-subtraction terms is exact (Section~\ref{sec:operator}).
    \item Phase covariance, amplitude homogeneity, locality, and equilibrium vanishing are proved (Section~\ref{sec:structure}).
    \item The equilibrium tangent limit is proved: linear Schr\"odinger dynamics is the linearisation of the full nonlinear dynamics at equilibrium (Section~\ref{sec:limit}).
    \item The quadratic example is computed explicitly (Section~\ref{sec:example}).
\end{itemize}

\subsection{Structural observations}

\begin{itemize}
    \item The closure nonlinearity belongs to the modulus-curvature class of nonlinear Schr\"odinger equations (Section~\ref{sec:NLS}).
    \item The coefficient $\sigma^2$ (normalised) and equilibrium subtraction $2\VW$ are closure-derived, not phenomenological (Section~\ref{sec:fixed}).
    \item The quantum potential of the Madelung formulation and the closure residual share the same functional form but differ in sign and role (Section~\ref{sec:madelung}).
\end{itemize}

\subsection{Open questions}

\begin{itemize}
    \item Does the closure nonlinear equation have a variational structure or a conserved energy functional?
    \item Can the nonlinear regime be constrained by existing experimental bounds on nonlinear quantum mechanics?
    \item Is the closure equation exactly equivalent to a specific Doebner--Goldin subfamily, or only structurally related?
    \item Does the nodal divergence of the residual indicate a physical boundary condition or a need for regularisation?
\end{itemize}

%==================================================================
\section{Conclusion}\label{sec:conclusion}
%==================================================================

Paper~XX places the closure residual into nonlinear quantum structure.

\textbf{Main result.} The closure-Nelson dynamics is a derived nonlinear Schr\"odinger equation of modulus-curvature type, with coefficient $\sigma^2$ (in Schr\"odinger-normalised form) and equilibrium subtraction $2\VW$ both fixed by the closure geometry. Linear quantum mechanics is not the full closure dynamics; it is the distinguished equilibrium tangent limit.

\textbf{What the closure framework provides.} Compared with a generic nonlinear Schr\"odinger model, the closure framework fixes: the linear Hamiltonian (Witten), the nonlinear functional form (modulus curvature), the nonlinear coefficient ($\sigma^2$ in normalised form), the equilibrium subtraction ($2\VW$), and the existence of a distinguished linear limit. These features are derived from the closure functional $\Fcal$ and the noise scale $\sigma$, not postulated.

\textbf{Programme arc.}
\begin{itemize}
    \item Papers~XII--XIV: dissipative closure dynamics,
    \item Paper~XV: phase and interference,
    \item Paper~XVI: evolution equation and structural gap,
    \item Paper~XVII: dissipative obstruction,
    \item Paper~XVIII: Nelson pairing and near-equilibrium Schr\"odinger recovery,
    \item Paper~XIX: isolation and classification of the closure residual,
    \item Paper~XX: structural placement as a derived nonlinear Schr\"odinger equation with equilibrium quantum limit.
\end{itemize}

The next stage is to determine whether the nonlinear residual is experimentally constrainable, whether it predicts observable deviations from standard quantum mechanics, and whether the closure framework thereby extends, modifies, or (in the equilibrium limit) exactly recovers the quantum-mechanical description of nature.

%==================================================================
\bibliographystyle{plainnat}
\bibliography{references}

\end{document}
